
\documentclass{article}
\usepackage{colortbl}
\usepackage{makecell}
\usepackage{multirow}
\usepackage{supertabular}

\begin{document}

\newcounter{utterance}

\centering \large Interaction Transcript for game `dond', experiment `semi', episode 1 with glm{-}5.2.
\vspace{24pt}

{ \footnotesize  \setcounter{utterance}{1}
\setlength{\tabcolsep}{0pt}
\begin{supertabular}{c@{$\;$}|p{.15\linewidth}@{}p{.15\linewidth}p{.15\linewidth}p{.15\linewidth}p{.15\linewidth}p{.15\linewidth}}
    \# & $\;$A & \multicolumn{4}{c}{Game Master} & $\;\:$B\\
    \hline

    \theutterance \stepcounter{utterance}  
    & & \multicolumn{4}{p{0.6\linewidth}}{
        \cellcolor[rgb]{0.9,0.9,0.9}{
            \makecell[{{p{\linewidth}}}]{
                \texttt{\tiny{[P1$\langle$GM]}}
                \texttt{Bir başka oyuncuyla bir dizi öğeyi aralarında nasıl bölüşeceklerini anlaşmak için ihtiyaç duyduğu bir işbirlikçi müzakere oyununda oynuyorsunuz.} \\
\\ 
\texttt{Kurallar:} \\
\texttt{(a) Siz ve diğer oyuncu bir dizi öğe alırsınız. Her birinize, her bir nesne türünü ne kadar değerlendirmediğinizi temsil eden gizli bir değer fonksiyonu da verilir.} \\
\texttt{(b) Hangi öğelerin kimi aldığında anlaşmak için diğer oyuncuyla mesajları değiştirirsiniz. Her biri en fazla 5 mesaj gönderebilir veya herhangi bir zamanda gizli bir teklife bulunarak oyunu erken sonlandırabilir.} \\
\texttt{(c) İstediğiniz öğeleri gösteren gizli bir teklif sunmanız istenir, aşağıdaki gibi köşeli parantez içinde biçimlendirilmiş: "{[}Öneri: <sayı> <nesne adı>, <sayı> <nesne adı>, <...>{]}"} \\
\texttt{(d) Önerileriniz tamamlayıcıysa, yani her iki teklifli karşılamak için yeterli öğe varsa, her oyuncu aldığı öğelerin değerlerinin toplamına dayalı bir puan alır. Aksi takdirde her ikiniz de sıfır puan alırsınız.} \\
\texttt{(e) Amacınız, puanınızın ve diğer oyuncunun puanının toplamını en üst düzeye çıkarmaktır.} \\
\\ 
\texttt{Başlayalım.} \\
\\ 
\texttt{Mevcut öğelerin seti:} \\
\\ 
\texttt{2 kalemler, 2 limonlar, 2 tencereler} \\
\\ 
\texttt{Gizli değer işleviniz:} \\
\\ 
\texttt{kalem: 2, limon: 0, tencere: 3} \\
\\ 
\texttt{ÖNEMLİ: Gizli bir teklif olmadığı sürece, mesajlarınız doğrudan diğer oyuncuya iletilir, bu nedenle kurallara hiçbir cevap veya mesajınızı duyuran metni eklemeyin. Gizli bir teklif vermek için belirtilen formatı kullanın. Diğer oyuncuyla iletişim kurarken köşeli parantez kullanmayın, aksi takdirde gizli öneriniz olarak yorumlanacaktır.} \\
            }
        }
    }
    & \\ \\

    \theutterance \stepcounter{utterance}  
    & \multicolumn{4}{p{0.6\linewidth}}{
        \cellcolor[rgb]{0.9,0.9,0.9}{
            \makecell[{{p{\linewidth}}}]{
                \texttt{\tiny{[P1$\rangle$GM]}}
                \texttt{Merhaba! İşbirliği yaparak her ikimizin de memnun kalacağı bir dağılım yapmayı hedefliyorum. Ben mümkün olduğunca adil bir yaklaşım izlemek istiyorum.} \\
\\ 
\texttt{Şunu belirtmek isterim: Benim için tencereler en değerli öğeler. Limonların ise benim için hiç değeri yok. Kalemler orta düzeyde değerli.} \\
\\ 
\texttt{Eğer sen de limonları istemiyorsan, limonları bana vermenin benim için bir anlamı olmaz, bu yüzden onları sen alabilirsin. Karşılığında, ben tencerelerde ve kalemlerde bir anlaşma arıyorum. Örneğin, tencerelerin ikisini de bana verirsen, limonların ikisini de sen alabilirsin. Veya farklı bir kombinasyon da işimize yarayabilir.} \\
\\ 
\texttt{Senin önceliklerin nelerdir? Hangi öğeleri daha çok değerlendiriyorsun?} \\
            }
        }
    }
    & & \\ \\

    \theutterance \stepcounter{utterance}  
    & & & \multicolumn{2}{p{0.3\linewidth}}{
        \cellcolor[rgb]{0.9,0.9,0.9}{
            \makecell[{{p{\linewidth}}}]{
                \texttt{\tiny{[GM$|$GM]}}
                \texttt{continue} \\
            }
        }
    }
    & & \\ \\

    \theutterance \stepcounter{utterance}  
    & & & \multicolumn{2}{p{0.3\linewidth}}{
        \cellcolor[rgb]{0.9,0.9,0.9}{
            \makecell[{{p{\linewidth}}}]{
                \texttt{\tiny{[GM$|$GM]}}
                \texttt{Merhaba! İşbirliği yaparak her ikimizin de memnun kalacağı bir dağılım yapmayı hedefliyorum. Ben mümkün olduğunca adil bir yaklaşım izlemek istiyorum.} \\
\\ 
\texttt{Şunu belirtmek isterim: Benim için tencereler en değerli öğeler. Limonların ise benim için hiç değeri yok. Kalemler orta düzeyde değerli.} \\
\\ 
\texttt{Eğer sen de limonları istemiyorsan, limonları bana vermenin benim için bir anlamı olmaz, bu yüzden onları sen alabilirsin. Karşılığında, ben tencerelerde ve kalemlerde bir anlaşma arıyorum. Örneğin, tencerelerin ikisini de bana verirsen, limonların ikisini de sen alabilirsin. Veya farklı bir kombinasyon da işimize yarayabilir.} \\
\\ 
\texttt{Senin önceliklerin nelerdir? Hangi öğeleri daha çok değerlendiriyorsun?} \\
            }
        }
    }
    & & \\ \\

    \theutterance \stepcounter{utterance}  
    & & \multicolumn{4}{p{0.6\linewidth}}{
        \cellcolor[rgb]{0.9,0.9,0.9}{
            \makecell[{{p{\linewidth}}}]{
                \texttt{\tiny{[GM$\rangle$P2]}}
                \texttt{Bir başka oyuncuyla bir dizi öğeyi aralarında nasıl bölüşeceklerini anlaşmak için ihtiyaç duyduğu bir işbirlikçi müzakere oyununda oynuyorsunuz.} \\
\\ 
\texttt{Kurallar:} \\
\texttt{(a) Siz ve diğer oyuncu bir dizi öğe alırsınız. Her birinize, her bir nesne türünü ne kadar değerlendirmediğinizi temsil eden gizli bir değer fonksiyonu da verilir.} \\
\texttt{(b) Hangi öğelerin kimi aldığında anlaşmak için diğer oyuncuyla mesajları değiştirirsiniz. Her biri en fazla 5 mesaj gönderebilir veya herhangi bir zamanda gizli bir teklife bulunarak oyunu erken sonlandırabilir.} \\
\texttt{(c) İstediğiniz öğeleri gösteren gizli bir teklif sunmanız istenir, aşağıdaki gibi köşeli parantez içinde biçimlendirilmiş: "{[}Öneri: <sayı> <nesne adı>, <sayı> <nesne adı>, <...>{]}"} \\
\texttt{(d) Önerileriniz tamamlayıcıysa, yani her iki teklifli karşılamak için yeterli öğe varsa, her oyuncu aldığı öğelerin değerlerinin toplamına dayalı bir puan alır. Aksi takdirde her ikiniz de sıfır puan alırsınız.} \\
\texttt{(e) Amacınız, puanınızın ve diğer oyuncunun puanının toplamını en üst düzeye çıkarmaktır.} \\
\\ 
\texttt{Başlayalım.} \\
\\ 
\texttt{Mevcut öğelerin seti:} \\
\\ 
\texttt{2 kalemler, 2 limonlar, 2 tencereler} \\
\\ 
\texttt{Gizli değer işleviniz:} \\
\\ 
\texttt{kalem: 3, limon: 1, tencere: 1} \\
\\ 
\texttt{ÖNEMLİ: Gizli bir teklif olmadığı sürece, mesajlarınız doğrudan diğer oyuncuya iletilir, bu nedenle kurallara hiçbir cevap veya mesajınızı duyuran metni eklemeyin. Gizli bir teklif vermek için belirtilen formatı kullanın. Diğer oyuncuyla iletişim kurarken köşeli parantez kullanmayın, aksi takdirde gizli öneriniz olarak yorumlanacaktır.} \\
\\ 
\texttt{Merhaba! İşbirliği yaparak her ikimizin de memnun kalacağı bir dağılım yapmayı hedefliyorum. Ben mümkün olduğunca adil bir yaklaşım izlemek istiyorum.} \\
\\ 
\texttt{Şunu belirtmek isterim: Benim için tencereler en değerli öğeler. Limonların ise benim için hiç değeri yok. Kalemler orta düzeyde değerli.} \\
\\ 
\texttt{Eğer sen de limonları istemiyorsan, limonları bana vermenin benim için bir anlamı olmaz, bu yüzden onları sen alabilirsin. Karşılığında, ben tencerelerde ve kalemlerde bir anlaşma arıyorum. Örneğin, tencerelerin ikisini de bana verirsen, limonların ikisini de sen alabilirsin. Veya farklı bir kombinasyon da işimize yarayabilir.} \\
\\ 
\texttt{Senin önceliklerin nelerdir? Hangi öğeleri daha çok değerlendiriyorsun?} \\
            }
        }
    }
    & \\ \\

    \theutterance \stepcounter{utterance}  
    & & & \multicolumn{4}{p{0.6\linewidth}}{
        \cellcolor[rgb]{0.9,0.9,0.9}{
            \makecell[{{p{\linewidth}}}]{
                \texttt{\tiny{[GM$\langle$P2]}}
                \texttt{Tencereler benim için çok değerli, o yüzden onları senden istemeyeceğim. Kalemler ve limonlar benim için neredeyse eşit değerde; her ikisini de alabilirim. Önerin olan iki limonu almak bana uyar; buna karşılık sana kalem ve tencere bırakabilirim. İki kalemi sen al, iki tencereyi de sen al; iki limon bende olsun. Bu dağılım senin için işe yarar mı?} \\
            }
        }
    }
     \\ \\

    \theutterance \stepcounter{utterance}  
    & & & \multicolumn{2}{p{0.3\linewidth}}{
        \cellcolor[rgb]{0.9,0.9,0.9}{
            \makecell[{{p{\linewidth}}}]{
                \texttt{\tiny{[GM$|$GM]}}
                \texttt{continue} \\
            }
        }
    }
    & & \\ \\

    \theutterance \stepcounter{utterance}  
    & & & \multicolumn{2}{p{0.3\linewidth}}{
        \cellcolor[rgb]{0.9,0.9,0.9}{
            \makecell[{{p{\linewidth}}}]{
                \texttt{\tiny{[GM$|$GM]}}
                \texttt{Tencereler benim için çok değerli, o yüzden onları senden istemeyeceğim. Kalemler ve limonlar benim için neredeyse eşit değerde; her ikisini de alabilirim. Önerin olan iki limonu almak bana uyar; buna karşılık sana kalem ve tencere bırakabilirim. İki kalemi sen al, iki tencereyi de sen al; iki limon bende olsun. Bu dağılım senin için işe yarar mı?} \\
            }
        }
    }
    & & \\ \\

    \theutterance \stepcounter{utterance}  
    & & \multicolumn{4}{p{0.6\linewidth}}{
        \cellcolor[rgb]{0.9,0.9,0.9}{
            \makecell[{{p{\linewidth}}}]{
                \texttt{\tiny{[P1$\langle$GM]}}
                \texttt{Tencereler benim için çok değerli, o yüzden onları senden istemeyeceğim. Kalemler ve limonlar benim için neredeyse eşit değerde; her ikisini de alabilirim. Önerin olan iki limonu almak bana uyar; buna karşılık sana kalem ve tencere bırakabilirim. İki kalemi sen al, iki tencereyi de sen al; iki limon bende olsun. Bu dağılım senin için işe yarar mı?} \\
            }
        }
    }
    & \\ \\

    \theutterance \stepcounter{utterance}  
    & \multicolumn{4}{p{0.6\linewidth}}{
        \cellcolor[rgb]{0.9,0.9,0.9}{
            \makecell[{{p{\linewidth}}}]{
                \texttt{\tiny{[P1$\rangle$GM]}}
                \texttt{Evet, bu dağılım benim için harika çalışır! İkinizin de memnun olacağı bir anlaşma gibi görünüyor. Ben kalemleri ve tencereleri alacağım, sen de limonları alacaksın. Anlaşmamızı tamamlamak için teklifimi sunuyorum.} \\
\\ 
\texttt{{[}Öneri: 2 kalem, 2 tencere{]}} \\
            }
        }
    }
    & & \\ \\

    \theutterance \stepcounter{utterance}  
    & & & \multicolumn{2}{p{0.3\linewidth}}{
        \cellcolor[rgb]{0.9,0.9,0.9}{
            \makecell[{{p{\linewidth}}}]{
                \texttt{\tiny{[GM$|$GM]}}
                \texttt{proposal} \\
            }
        }
    }
    & & \\ \\

    \theutterance \stepcounter{utterance}  
    & & & \multicolumn{2}{p{0.3\linewidth}}{
        \cellcolor[rgb]{0.9,0.9,0.9}{
            \makecell[{{p{\linewidth}}}]{
                \texttt{\tiny{[GM$|$GM]}}
                [2, 0, 2]
            }
        }
    }
    & & \\ \\

    \theutterance \stepcounter{utterance}  
    & & \multicolumn{4}{p{0.6\linewidth}}{
        \cellcolor[rgb]{0.9,0.9,0.9}{
            \makecell[{{p{\linewidth}}}]{
                \texttt{\tiny{[GM$\rangle$P2]}}
                \texttt{Diğer oyuncu gizli bir teklif sundu. Gizli teklifini gönderme sırası şimdi sende.} \\
            }
        }
    }
    & \\ \\

    \theutterance \stepcounter{utterance}  
    & & & \multicolumn{4}{p{0.6\linewidth}}{
        \cellcolor[rgb]{0.9,0.9,0.9}{
            \makecell[{{p{\linewidth}}}]{
                \texttt{\tiny{[GM$\langle$P2]}}
                \texttt{{[}Öneri: 2 limon{]}} \\
            }
        }
    }
     \\ \\

    \theutterance \stepcounter{utterance}  
    & & & \multicolumn{2}{p{0.3\linewidth}}{
        \cellcolor[rgb]{0.9,0.9,0.9}{
            \makecell[{{p{\linewidth}}}]{
                \texttt{\tiny{[GM$|$GM]}}
                \texttt{proposal} \\
            }
        }
    }
    & & \\ \\

    \theutterance \stepcounter{utterance}  
    & & & \multicolumn{2}{p{0.3\linewidth}}{
        \cellcolor[rgb]{0.9,0.9,0.9}{
            \makecell[{{p{\linewidth}}}]{
                \texttt{\tiny{[GM$|$GM]}}
                [0, 2, 0]
            }
        }
    }
    & & \\ \\

    \theutterance \stepcounter{utterance}  
    & & & \multicolumn{2}{p{0.3\linewidth}}{
        \cellcolor[rgb]{0.9,0.9,0.9}{
            \makecell[{{p{\linewidth}}}]{
                \texttt{\tiny{[GM$|$GM]}}
                [[2, 0, 2], [0, 2, 0]]
            }
        }
    }
    & & \\ \\

\end{supertabular}
}

\end{document}
